% Reusable research-strand triangle.
% Usage: \ResearchTriangle{mode}
%        \ResearchTriangleCorner{mode}
%   mode ∈ {off, none, all, stats-physics, stats-fp, fp-physics}
%
% Vertices (clockwise from top):
%   Statistics
%   Functional Programming
%   Physics
%
% Edges:
%   stats-physics : Statistics -- Physics              (HMC / spin-fermion)
%   stats-fp      : Statistics -- Functional Programming (FRP)
%   fp-physics    : Functional Programming -- Physics  (typed fusion)

\newcommand{\currentEdge}{off}

\newcommand{\ResearchTriangle}[1]{%
  \begin{tikzpicture}[
    baseline=(current bounding box.center),
    every node/.style={font=\sffamily},
    edge dim/.style={draw=black!28, line width=1.0pt, line cap=round},
    edge hi/.style={draw=black!88, line width=2.4pt, line cap=round},
    vtx/.style={circle, fill=black!88, draw=white, line width=1.2pt, inner sep=0pt, minimum size=2.4mm},
    lab/.style={align=center, font=\small\sffamily\bfseries, text=black!88, inner sep=1pt},
  ]
    \def\triMode{#1}
    \coordinate (S) at (90:2.35);
    \coordinate (F) at (210:2.35);
    \coordinate (P) at (330:2.35);

    \draw[edge dim] (S) -- (F);
    \draw[edge dim] (S) -- (P);
    \draw[edge dim] (F) -- (P);

    \ifstrequal{\triMode}{all}{%
      \draw[edge hi] (S) -- (F);
      \draw[edge hi] (S) -- (P);
      \draw[edge hi] (F) -- (P);
    }{}
    \ifstrequal{\triMode}{stats-fp}{%
      \draw[edge hi] (S) -- (F);
    }{}
    \ifstrequal{\triMode}{bayes-fp}{%
      \draw[edge hi] (S) -- (F);
    }{}
    \ifstrequal{\triMode}{stats-physics}{%
      \draw[edge hi] (S) -- (P);
    }{}
    \ifstrequal{\triMode}{bayes-physics}{%
      \draw[edge hi] (S) -- (P);
    }{}
    \ifstrequal{\triMode}{fp-physics}{%
      \draw[edge hi] (F) -- (P);
    }{}

    \node[vtx] at (S) {};
    \node[vtx] at (F) {};
    \node[vtx] at (P) {};
    \node[lab, above=10pt] at (S) {Statistics};
    \node[lab, below=10pt] at (F) {Functional\\Programming};
    \node[lab, below=10pt] at (P) {Physics};
  \end{tikzpicture}%
}

\newcommand{\ResearchTriangleCorner}[1]{%
  \ifstrequal{#1}{off}{}{%
    \begin{tikzpicture}[
      baseline=(current bounding box.center),
      every node/.style={font=\tiny\sffamily\bfseries},
      edge dim/.style={draw=black!22, line width=0.55pt, line cap=round},
      edge hi/.style={draw=black!88, line width=1.35pt, line cap=round},
      vtx/.style={circle, fill=black!88, draw=white, line width=0.7pt, inner sep=0pt, minimum size=1.5mm},
    ]
      \def\triMode{#1}
      \coordinate (S) at (90:0.78);
      \coordinate (F) at (210:0.78);
      \coordinate (P) at (330:0.78);

      \draw[edge dim] (S) -- (F);
      \draw[edge dim] (S) -- (P);
      \draw[edge dim] (F) -- (P);

      \ifstrequal{\triMode}{all}{%
        \draw[edge hi] (S) -- (F);
        \draw[edge hi] (S) -- (P);
        \draw[edge hi] (F) -- (P);
      }{}
      \ifstrequal{\triMode}{stats-fp}{%
        \draw[edge hi] (S) -- (F);
      }{}
      \ifstrequal{\triMode}{bayes-fp}{%
        \draw[edge hi] (S) -- (F);
      }{}
      \ifstrequal{\triMode}{stats-physics}{%
        \draw[edge hi] (S) -- (P);
      }{}
      \ifstrequal{\triMode}{bayes-physics}{%
        \draw[edge hi] (S) -- (P);
      }{}
      \ifstrequal{\triMode}{fp-physics}{%
        \draw[edge hi] (F) -- (P);
      }{}

      \node[vtx] at (S) {};
      \node[vtx] at (F) {};
      \node[vtx] at (P) {};
      \node[above=3pt] at (S) {S};
      \node[below=3pt] at (F) {F};
      \node[below=3pt] at (P) {P};
    \end{tikzpicture}%
  }%
}
